\section{Literature Review}
Short-term load forecasting (STLF) and regional grid risk management traditionally relied on linear statistical time-series models such as AutoRegressive Integrated Moving Average (ARIMA) and classical statistical principles \cite{hyndman2018forecasting}. While effective for stationary univariate series, statistical methods struggle to capture nonlinear interactions between environmental parameters and grid utilization dynamics. Inductive machine learning algorithms \cite{michalski1983theory}, including Random Forests \cite{breiman2001random} and XGBoost \cite{chen2016xgboost}, significantly improved tabular forecasting accuracy by modeling high-dimensional feature interactions, but they treat telemetry records as independent samples and discard chronological temporal structures.

To model temporal dependencies, recurrent neural architectures—specifically Long Short-Term Memory (LSTM) \cite{hochreiter1997long} and Gated Recurrent Unit (GRU) \cite{cho2014learning, chung2014empirical} networks—have been widely adopted in power systems engineering to overcome vanishing gradient challenges in deep sequences \cite{bengio1994learning}. Hybrid variants, such as CNN-LSTM models \cite{shi2015convolutional} and temporal convolution models \cite{lai2018modeling}, use 1D convolutional layers to extract localized feature patterns before passing hidden representations to recurrent units.

Recently, time-series Transformers have emerged as a prominent paradigm for long sequence modeling. Self-attention mechanisms \cite{vaswani2017attention} enable direct temporal modeling. Architectures such as Informer \cite{zhou2021informer} utilize ProbSparse self-attention to reduce memory and time complexity from $O(L^2)$ to $O(L \log L)$. Autoformer \cite{wu2021autoformer} introduces auto-correlation mechanisms, PatchTST \cite{nie2022time} processes sub-series patches to preserve local semantics, Temporal Fusion Transformers (TFT) \cite{lim2021temporal} combine gating mechanisms with multi-head attention, and TimesNet \cite{wu2023timesnet} transforms 1D series into 2D spaces to model intra-period and inter-period variations. However, recent empirical audits by Zeng et al. \cite{zeng2023transformers} suggest that complex self-attention layers do not consistently outperform simpler linear baselines across all time-series tasks, highlighting the necessity of empirical benchmarking on real-world grid telemetry.

In the context of developing power grids, national capacity deficits and infrastructure challenges in Bangladesh \cite{sheikh2019power, bpdb2023report} create frequent unannounced load-shedding events. Prior smart meter load management studies \cite{rashid2016automated} and machine learning grid monitoring frameworks \cite{lipu2023intelligent} demonstrate the value of automated telemetry. Meteorological parameters ingested via weather APIs \cite{openmeteo2024} provide critical environmental context. However, existing literature rarely couples localized Upazila-level sequence predictions with explainable artificial intelligence (XAI) frameworks \cite{adadi2018peeking} to generate automated, rule-based natural language alerts for grid operators. Table \ref{tab:related_works} summarizes key related literature in relation to our proposed framework.

\begin{table}[htbp]
\caption{Summary of Related Works and Proposed Framework}
\centering
\resizebox{\columnwidth}{!}{%
\begin{tabular}{lp{1.8cm}p{2.2cm}p{2.5cm}}
\toprule
\textbf{Study} & \textbf{Methodology} & \textbf{Target Dataset} & \textbf{Key Limitation / Gap} \\
\midrule
Rashid et al. \cite{rashid2016automated} & Smart Meters & Regional Grid Telemetry & Lacks deep sequence modeling \\
Sheikh et al. \cite{sheikh2019power} & Statistical Analysis & Bangladesh Grid & Lacks predictive forecasting \\
Lipu et al. \cite{lipu2023intelligent} & Machine Learning & Grid Telemetry & No real-time explainability layer \\
Zhou et al. \cite{zhou2021informer} & Informer & Energy Benchmarks & Lacks operational alert integration \\
Wu et al. \cite{wu2023timesnet} & TimesNet & General Time-Series & High training time complexity \\
\textbf{Proposed Framework} & \textbf{Informer + Rule Engine} & \textbf{Bangladesh Grid (143 Upazilas)} & \textbf{Evaluated on two distribution divisions} \\
\bottomrule
\end{tabular}%
}
\label{tab:related_works}
\end{table}
